We have made many conjectures.  We hope that our conjectures are true, but some are likely false. How do we justify that a conjecture is true or false?  When a conjecture is true, we write a proof.  When a conjecture is false, we usually provide an example showing that it is false (a \vocab{counterexample}).  In deciding whether a conjecture is true or false,  developing a proof, and building counterexamples, we rely on logic. 

Many exercises in this textbook require an explanation or justification. As computers are capable of more and more computations, it is the explanation, the human part, that becomes increasingly valuable to learn.  It can be difficult to know how much you need to say for an explanation or proof.  Understanding the logic of the statement can help. Logic is also the backbone of how computers, computer programming, databases, and search engines work.  

This chapter introduces formal logic, including \Cref{sec_logic_connectives} {\seclogicconn}, \Cref{sec_quant_cex} {\secquant}, \Cref{sec_conditional_statements} {\secifthen}, and \Cref{sec_converse_cp} {\secconvcp}.  Throughout these sections, we construct counterexamples, learn many new proof formats, and explore logical equivalence.  We introduce logical inference in \Cref{sec_inference} {\secinference}.




